\subsection{Solucionario: Cálculo}

% ---  1 ---
\begin{sol}[{Cálculo Unidimensional}]
\begin{enumerate}
    \item $\frac{d}{dt} [t^2 \cos(3t)]$: Usando la regla del producto $(uv)' = u'v + uv'$:
    $2t \cos(3t) + t^2 (-\sin(3t) \cdot 3) = 2t\cos(3t) - 3t^2\sin(3t)$.
    
    \item $\frac{d}{dt} [ \ln(t^2 + 1) e^{-2t} ]$: Regla del producto y cadena:
    $\frac{2t}{t^2+1} e^{-2t} + \ln(t^2+1) (-2e^{-2t}) = e^{-2t} \left( \frac{2t}{t^2+1} - 2\ln(t^2+1) \right)$.
    
    \item $\int (3t^2 - 5\sin(t)) dt$: Integración directa:
    $t^3 - 5(-\cos(t)) + C = t^3 + 5\cos(t) + C$.
    
    \item $\int x e^{x^2} dx$: Sustitución $u = x^2$, $du = 2x dx$:
    $\frac{1}{2} \int e^u du = \frac{1}{2} e^{x^2} + C$.
\end{enumerate}
\end{sol}

% ---  2 ---
\begin{sol}[{ Cálculo Multivariable}]
\begin{enumerate}
    \item $\frac{\partial}{\partial x}(x^3 y^2 + \sin(xy))$: Se trata a $y$ como constante:
    $3x^2 y^2 + \cos(xy) \cdot y = 3x^2 y^2 + y\cos(xy)$.
    
    \item $\frac{\partial^2}{\partial x \partial y} (e^{x^2+y^2})$: Primero derivamos respecto a $y$: $\frac{\partial}{\partial y} = 2y e^{x^2+y^2}$. Luego respecto a $x$: $2y (e^{x^2+y^2} \cdot 2x) = 4xy e^{x^2+y^2}$.
    
    \item $\int_0^1 \int_0^2 (x + 2y) \, dy \, dx$: 
    $\int_0^1 [xy + y^2]_0^2 dx = \int_0^1 (2x + 4) dx = [x^2 + 4x]_0^1 = 1 + 4 = 5$.
    
    \item $\int_0^\pi \int_0^R r \, dr \, d\theta$: 
    $\int_0^\pi [\frac{r^2}{2}]_0^R d\theta = \frac{R^2}{2} \int_0^\pi d\theta = \frac{\pi R^2}{2}$.
\end{enumerate}
\end{sol}

% ---  3 ---
\begin{sol}[{ Derivación Vectorial}]
\begin{enumerate}
    \item $\frac{d}{dt} (\vec{u} \cdot \vec{v}) = \dot{\vec{u}} \cdot \vec{v} + \vec{u} \cdot \dot{\vec{v}}$.
    \item $\vec{r}(t) = \cos(\omega t)\vec{i} + \sin(\omega t)\vec{j} + bt\vec{k}$. 
    Derivando componente a componente:
    $\dot{\vec{r}} = -\omega\sin(\omega t)\vec{i} + \omega\cos(\omega t)\vec{j} + b\vec{k}$.
    $\ddot{\vec{r}} = -\omega^2\cos(\omega t)\vec{i} - \omega^2\sin(\omega t)\vec{j}$.
    \item $\frac{d}{dt} (\frac{1}{2} m \vec{v} \cdot \vec{v}) = \frac{1}{2} m (\dot{\vec{v}} \cdot \vec{v} + \vec{v} \cdot \dot{\vec{v}}) = m \vec{v} \cdot \dot{\vec{v}} = m \vec{v} \cdot \vec{a} = P$ (Potencia).
\end{enumerate}
\end{sol}

% ---  4 ---
\begin{sol}[{ Bases Intrínsecas}]
Para $\vec{r}(s) = \left( a \cos(\frac{s}{c}), a \sin(\frac{s}{c}), \frac{bs}{c} \right)$ con $c = \sqrt{a^2+b^2}$:
\begin{enumerate}
    \item $\hat{T} = \frac{d\vec{r}}{ds} = \frac{1}{c} \left( -a\sin(\frac{s}{c}), a\cos(\frac{s}{c}), b \right)$.
    \item $\frac{d\hat{T}}{ds} = \frac{1}{c^2} (-a\cos(\frac{s}{c}), -a\sin(\frac{s}{c}), 0)$. La norma es $\frac{a}{c^2}$.
    $\hat{N} = \frac{d\hat{T}/ds}{a/c^2} = (-\cos(\frac{s}{c}), -\sin(\frac{s}{c}), 0)$.
    \item $\hat{B} = \hat{T} \times \hat{N} = \frac{1}{c} (b\sin(\frac{s}{c}), -b\cos(\frac{s}{c}), a)$.
\end{enumerate}
\end{sol}

% ---  5 ---
\begin{sol}[{ Integrales Triples}]
\begin{enumerate}
    \item $J = \rho \int_0^H \int_0^{2\pi} \int_0^R r^2 (r \, dr \, d\theta \, dz)$: 
    $J = \rho \cdot H \cdot 2\pi \cdot \frac{R^4}{4} = \frac{1}{2} (\rho \pi R^2 H) R^2 = \frac{1}{2} M R^2$.
    \item $P_{xy} = \rho \int_0^L x \, dx \int_0^L y \, dy \int_0^L dz = \rho (\frac{L^2}{2})(\frac{L^2}{2})(L) = \frac{1}{4} \rho L^5 = \frac{1}{4} ML^2$.
\end{enumerate}
\end{sol}

% ---  6 ---
\begin{sol}[{ Centro de Masa}]
\begin{enumerate}
    \item $\vec{r}_{cm} = \frac{\sum m_i \vec{r}_i}{\sum m_i} = \frac{2(1,0,0) + 3(0,2,0) + 5(0,0,2)}{2+3+5} = \frac{(2, 6, 10)}{10} = (0.2, 0.6, 1.0)$.
    \item $\bar{y} = \frac{1}{M} \iint y \rho dA$. En polares: $y = r\sin\theta$, $dA = r \, dr \, d\theta$.
    $M = \frac{1}{2}\pi R^2 \rho$. 
    $\iint y \rho dA = \rho \int_0^\pi \int_0^R r^2 \sin\theta \, dr \, d\theta = \rho \frac{R^3}{3} [-\cos\theta]_0^\pi = \frac{2\rho R^3}{3}$.
    $\bar{y} = \frac{2\rho R^3 / 3}{\pi R^2 \rho / 2} = \frac{4R}{3\pi}$.
    \item $M = \int_0^L kx \, dx = \frac{1}{2}kL^2$. 
    $\bar{x} = \frac{1}{M} \int_0^L x (kx) dx = \frac{1}{M} \frac{kL^3}{3} = \frac{kL^3/3}{kL^2/2} = \frac{2}{3}L$.
\end{enumerate}
\end{sol}